% © 2026 Ing. David Jaroš — CC BY-NC-ND 4.0
%
% This work is licensed under a Creative Commons Attribution-NonCommercial-NoDerivatives
% 4.0 International License (CC BY-NC-ND 4.0).
%
% License History: Earlier drafts (up to v0.3) were released under CC BY 4.0.
% From v0.4 onward, all material is released under CC BY-NC-ND 4.0 to protect
% the integrity of the theoretical work during ongoing academic development.
%
% Three Generations Research Track — Jacobi Theta Step 1
% Status: Working document — conclusions labelled CONJECTURE or PROVED
% Version: 1.0
% Date: 2026-03-05

\documentclass[12pt,a4paper]{article}

\usepackage{amsmath,amssymb,amsthm}
\usepackage{hyperref}
\usepackage{geometry}
\usepackage{enumitem}
\geometry{margin=2.5cm}

\newtheorem{theorem}{Theorem}[section]
\newtheorem{lemma}[theorem]{Lemma}
\newtheorem{proposition}[theorem]{Proposition}
\newtheorem{corollary}[theorem]{Corollary}
\newtheorem{definition}[theorem]{Definition}
\theoremstyle{remark}
\newtheorem{remark}[theorem]{Remark}
\newtheorem{conjecture}[theorem]{Conjecture}

\newcommand{\C}{\mathbb{C}}
\newcommand{\R}{\mathbb{R}}
\newcommand{\H}{\mathbb{H}}
\newcommand{\Z}{\mathbb{Z}}
\newcommand{\Sc}{\mathrm{Sc}}
\newcommand{\varthree}{\vartheta_3}

\title{%
  Jacobi Theta Functions as Generation Index in UBT\\
  Step 1: Fourier vs.\ Taylor Mode Expansion\\
  \large Three Generations Research Track --- Jacobi Theta Step 1
}

\author{David Jaro\v{s}}
\date{2026-03-05}

\begin{document}

\maketitle

\begin{abstract}
The UBT field $\Theta(x,\tau)$ with $\tau = t + i\psi$ and $\psi$ compact
($\psi \sim \psi + 2\pi R_\psi$) admits two natural mode expansions: a
Taylor series in $\psi$ and a Fourier series in $\psi$.  We derive the
precise relation between the Taylor modes $\Theta_n$ (used in
\texttt{st3\_complex\_time\_generations.tex}) and the Fourier modes
$\hat{\Theta}_n$ (the natural basis for a compact direction), and explain
why the Fourier basis is more natural geometrically.  We then compare the
$n^2$ Jacobi theta spectrum with the experimental lepton-mass ratios and
document the mismatch precisely, which specifies what additional dynamics
are needed.
\end{abstract}

\tableofcontents

%=============================================================
\section{Setup: Compact Imaginary Time and Two Mode Expansions}
\label{sec:setup}
%=============================================================

Throughout this document we adopt the conventions of
\texttt{st3\_complex\_time\_generations.tex}: complex time is
$\tau = t + i\psi$, and $\psi$ is compact with period $2\pi R_\psi$:
\begin{equation}
  \psi \sim \psi + 2\pi R_\psi .
\end{equation}
The field $\Theta(x,\tau)$ is valued in $\C\otimes\H$ (biquaternions)
and is smooth in $\psi$.

\subsection{Taylor expansion (used in ST-3)}

Since $\Theta$ is smooth in $\psi$ we may write
\begin{equation}
  \label{eq:taylor}
  \Theta(x,t,\psi)
  = \sum_{n=0}^{\infty} \frac{\psi^n}{n!}\,\Theta_n(x,t),
  \qquad
  \Theta_n(x,t) \;=\; \left.\frac{\partial^n\Theta}{\partial\psi^n}
                       \right|_{\psi=0}.
\end{equation}
This is the expansion used in ST-3 to define the
\emph{Taylor modes} $\Theta_n$.

\subsection{Fourier expansion (natural for compact $\psi$)}

Because $\psi$ is compact the more natural decomposition is the Fourier
series
\begin{equation}
  \label{eq:fourier}
  \Theta(x,t,\psi)
  = \sum_{n=-\infty}^{\infty}
    \hat{\Theta}_n(x,t)\,e^{in\psi/R_\psi},
\end{equation}
with inverse
\begin{equation}
  \label{eq:fourier_inv}
  \hat{\Theta}_n(x,t)
  = \frac{1}{2\pi R_\psi}
    \int_0^{2\pi R_\psi} \Theta(x,t,\psi)\,e^{-in\psi/R_\psi}\,d\psi.
\end{equation}

%=============================================================
\section{Relation Between Taylor and Fourier Modes}
\label{sec:relation}
%=============================================================

\begin{proposition}[Taylor--Fourier relation]
\label{prop:tf_relation}
Denote the Fourier mode vector as $\hat\Theta_n$ and the Taylor mode
as $\Theta_n = \partial_\psi^n\Theta|_{\psi=0}$.  Then
\begin{equation}
  \label{eq:tf_relation}
  \Theta_n(x,t)
  = \sum_{k=-\infty}^{\infty}
    \left(\frac{ik}{R_\psi}\right)^n \hat\Theta_k(x,t).
\end{equation}
\end{proposition}

\begin{proof}
Differentiate the Fourier series \eqref{eq:fourier} $n$ times with
respect to $\psi$:
\[
  \frac{\partial^n\Theta}{\partial\psi^n}
  = \sum_k \hat\Theta_k(x,t)
    \left(\frac{ik}{R_\psi}\right)^n e^{ik\psi/R_\psi}.
\]
Setting $\psi=0$ gives \eqref{eq:tf_relation}.
\end{proof}

\begin{remark}
The formula \eqref{eq:tf_relation} shows that every Taylor mode
$\Theta_n$ is an infinite sum of \emph{all} Fourier modes.  Conversely,
a single Fourier mode $\hat\Theta_k$ contributes to every Taylor mode.
The two bases are \emph{not} in one-to-one correspondence unless $\Theta$
contains only a single Fourier harmonic.
\end{remark}

\subsection{Which basis is more natural?}

For a \textbf{compact} direction $\psi$, the Fourier basis is more natural
for the following reasons:
\begin{enumerate}[label=(\roman*)]
  \item \textbf{Completeness and unitarity:}  The functions
        $\{e^{in\psi/R_\psi}\}_{n\in\Z}$ form an orthonormal basis of
        $L^2([0,2\pi R_\psi])$.  The Taylor polynomials $\psi^n/n!$ do not
        form a basis of the same space.
  \item \textbf{Eigenmode structure:}  Each Fourier mode $\hat\Theta_n$ is
        an eigenfunction of the $\psi$-Laplacian $\partial_\psi^2$ with
        eigenvalue $-(n/R_\psi)^2$.  Taylor modes are not eigenfunctions.
  \item \textbf{Kaluza--Klein analogy:}  In standard Kaluza--Klein
        dimensional reduction the physical spectrum is labelled by the
        Fourier mode number $n$, not by the order of a Taylor coefficient.
  \item \textbf{Periodicity:}  Taylor modes in general do not respect the
        compactness condition $\Theta(\psi+2\pi R_\psi)=\Theta(\psi)$;
        Fourier modes do by construction.
\end{enumerate}

\begin{remark}[Taylor modes as a practical tool]
Despite the geometric superiority of the Fourier basis, the Taylor modes
are a useful truncation tool near $\psi=0$.  The results of ST-3
(independence, same SU(3) quantum numbers, $\psi$-parity) hold in both
bases, since they follow from the algebraic structure of $\C\otimes\H$
and the symmetry $\psi\to-\psi$, not from the specific expansion basis.
\end{remark}

%=============================================================
\section{Jacobi Theta Structure and the $n^2$ Spectrum}
\label{sec:jacobi}
%=============================================================

\subsection{The Jacobi theta function}

The Jacobi theta function of the third kind is
\begin{equation}
  \label{eq:jacobi_theta3}
  \varthree(z,\tau)
  = \sum_{n=-\infty}^{\infty} q^{n^2}\,e^{2\pi inz},
  \qquad q = e^{i\pi\tau}.
\end{equation}
With $\tau = t + i\psi$ and $q = e^{i\pi(t+i\psi)} = e^{i\pi t-\pi\psi}$
the Fourier coefficients are
\begin{equation}
  \label{eq:jacobi_coeff}
  c_n(\tau) = q^{n^2} = e^{i\pi n^2 t}\,e^{-\pi n^2\psi}.
\end{equation}

\subsection{Comparison with the UBT Fourier expansion}

In the UBT Fourier expansion \eqref{eq:fourier} the Fourier coefficient
of mode $n$ is $\hat\Theta_n(x,t)$ — an arbitrary $\C\otimes\H$-valued
function.  For the expansion to coincide \emph{structurally} with
\eqref{eq:jacobi_theta3}, the coefficient envelope must satisfy
\begin{equation}
  \label{eq:jacobi_constraint}
  |\hat\Theta_n| \;\propto\; q^{n^2}
  \quad\Longleftrightarrow\quad
  |\hat\Theta_n| \;\propto\; e^{-\pi n^2\psi_0}
\end{equation}
for some reference value $\psi_0 > 0$.

\textbf{Open question:}  Does the UBT action $S[\Theta]$ select a
saddle-point configuration in which the envelope of the Fourier
coefficients satisfies \eqref{eq:jacobi_constraint}?

A partial answer is given in Appendix~H (\texttt{Appendix\_H\_Theta\_Phase\_Emergence.tex}): the drift-diffusion steady-state
solutions are Jacobi theta functions.  This supports the claim but a
rigorous derivation from $S[\Theta]$ remains \textbf{[OPEN]}.

%=============================================================
\section{Mass Hierarchy from the $n^2$ Spectrum}
\label{sec:mass_n2}
%=============================================================

If the Fourier coefficients follow $|\hat\Theta_n|\propto e^{-\pi n^2\psi_0}$
then the effective squared mass of mode $n$ relative to mode $0$ is
controlled by the ratio of norms, giving a spectrum
\begin{equation}
  \label{eq:n2_spectrum}
  m_n^2 \;\propto\; n^2
  \quad\text{(from the Jacobi $n^2$ exponent)}.
\end{equation}

\subsection{Comparison with experiment}

The three charged lepton masses are (PDG values, $\text{MeV}/c^2$):
\begin{align}
  m_e    &\approx 0.511  &\text{(electron, generation 0)}, \\
  m_\mu  &\approx 105.7  &\text{(muon, generation 1)}, \\
  m_\tau &\approx 1776.9 &\text{(tau, generation 2)}.
\end{align}
The ratios are
\begin{equation}
  m_e : m_\mu : m_\tau
  \;\approx\; 1 : 207 : 3477.
\end{equation}

The $n^2$ Jacobi spectrum predicts (with $m_0\equiv 1$):
\begin{equation}
  m_0 : m_1 : m_2
  \;=\; 0^2 : 1^2 : 2^2
  \;=\; 0 : 1 : 4.
\end{equation}

\subsection{The mismatch}

\begin{enumerate}[label=(\roman*)]
  \item The $n^2$ spectrum predicts $m_1/m_0 = 1$ and $m_2/m_0 = 4$.
        Experiment gives $207$ and $3477$.  The mismatch is a factor of
        $\approx 207$ at the muon level and $\approx 870$ at the tau level.
  \item The \emph{ratio} $m_2/m_1$ is particularly constraining:
        the $n^2$ spectrum predicts $4/1=4$, while the experiment gives
        $3477/207 \approx 16.8$.
  \item The experimental ratio $m_2/m_1 \approx 16.8 \approx 4^{2.03}$
        is suggestive of a \emph{quartic} (rather than quadratic) scaling
        in the mode number, but this is only a rough guide.
  \item An exponential fit $m_n \propto e^{bn}$ with $b=\ln 207\approx
        5.33$ gives $m_0:m_1:m_2 = 1:207:42849$, which overshoots
        $m_\tau$ by a factor $\approx 12.3$.
\end{enumerate}

\begin{remark}[Documenting the mismatch]
The $n^2$ Jacobi spectrum is geometrically natural (it arises directly
from the modular parameter $q^{n^2}$) but is inconsistent with the
experimental mass ratios by more than two orders of magnitude.  This
mismatch is not a failure of the approach; it tells us precisely what
\emph{additional} dynamics are required beyond the free Fourier structure.
Section~3 of \texttt{step3\_mass\_from\_psi\_energy.tex} characterises
the required interaction $V(\Theta)$.
\end{remark}

\subsection{Summary of Step 1}

\begin{itemize}
  \item[\checkmark] Precise relation between Taylor and Fourier modes:
        equation~\eqref{eq:tf_relation} --- \textbf{PROVED}.
  \item[\checkmark] Fourier modes are the natural basis for compact
        $\psi$ --- \textbf{ARGUED} (four reasons, Section~\ref{sec:relation}).
  \item[\checkmark] $n^2$ Jacobi spectrum vs.\ experimental ratios:
        mismatch precisely documented --- see Section~\ref{sec:mass_n2}.
  \item[$\square$] Whether $S[\Theta]$ selects $q^{n^2}$ coefficients:
        \textbf{[OPEN]} --- partial evidence from Appendix~H.
\end{itemize}

\end{document}
